\section{History: The Documented Development and the Received Tradition}
\label{sec:history}

The Rosary's history must be told twice, and honestly both times. There
is a devotional and papal tradition, repeated for centuries in the
Breviary lessons and in papal documents, that the Blessed Virgin gave
the Rosary to St.~Dominic (d.~1221) as a weapon against the Albigensian
heresy. And there is the documented development, traced by Catholic
scholarship since the eighteenth century and stated soberly by the
modern magisterium itself: the Rosary ``gradually took form in the
second millennium'' (RVM 1). The two accounts are treated in turn; the
governing rule is Herbert Thurston's own---``the freest criticism of the
historical origin of the devotion, which involves no point of doctrine,
is compatible with a full appreciation of the devotional treasures which
this pious exercise brings within the reach of all.''\endnote{Herbert
Thurston and Andrew Shipman, ``The Rosary,'' \work{The Catholic
Encyclopedia}, vol.~13 (New York: Robert Appleton, 1912),
\url{https://www.newadvent.org/cathen/13184b.htm}. Except where another
witness is named, the documented stages in this section, with their
medieval sources, follow this article, which remains the classic
English summary of the evidence assembled by the Bollandists, T.~Esser,
and the scholarship of Thurston's generation; its findings are reported
as checked scholarship, not re-verified in the medieval manuscripts.}

\subsection{Counted prayer before the Rosary}

Nothing in the Rosary is older than the impulse to count prayers.
Thurston's survey ranges deliberately wide---an Assyrian relief at
Nineveh whose winged figures hold what its excavator called ``a garland
or rosary''; the Muslim \term{tasbih} of thirty-three, sixty-six, or
ninety-nine beads; the king of Malabar whose hundred-odd precious
stones astonished Marco Polo; the Buddhist strings that astonished
Francis Xavier in Japan---not to derive the Rosary from any of them, but
to establish that a counting string is a natural human device and proves
nothing by itself about who invented what.\endnote{Thurston, ``The
Rosary,'' opening survey. The comparative material establishes the
banality of prayer-counting; it is not offered here, and is not offered
there, as evidence of borrowing.}

Within Christianity the practice is monastic and very early. Palladius
records Paul the Hermit, in the fourth century, laying up three hundred
pebbles and throwing one away at each of the three hundred prayers he
had imposed on himself; beside the body of a fourth-century ascetic
disinterred at Antino\"e was found a pierced board that has generally
been read as a prayer-counter; and a privilege of Hadrian~I in 782
required the monks of St.~Apollinaris in Classe to say \term{Kyrie
eleison} three hundred times twice a day, a burden that all but
presupposes some apparatus.\endnote{Thurston, ``The Rosary,'' citing
Palladius, \work{Historia Lausiaca} xx (with Butler II, 63); the
Antino\"e find; and Jaff\'e--L\"owenfeld n.~2437 for Hadrian~I's
privilege of 782. Loci are reported from Thurston, not re-verified in
the editions named.} In the Christian East the counted cord never
became the Rosary: it remained the monastic \term{komboschoinion} or
\term{kombologion}, conferred with the full habit as the monk's
``spiritual sword'' and told with the Jesus Prayer or its penitential
form, under the rule that the Nomocanon states flatly---``The rosary
should have one hundred beads; and upon each bead the prescribed prayer
should be recited.'' The Western Rosary has not become general among the
Eastern faithful, and where it appears among Eastern Catholic laity it
is an imitation of Roman practice rather than an indigenous
growth.\endnote{Thurston and Shipman, ``The Rosary,'' section ``In the
Greek Church'' (Shipman's contribution): the Greek \term{kombologion}
and \term{komboschoinion}, the Russian \term{vervitza}, \term{chotki},
and \term{liestovka}, the Romanian \term{matanie}; Nomocanon rule 87;
the invocation ``O Lord Jesus Christ, Son and Word of the living God,
through the intercessions of thy immaculate Mother and of all thy
Saints, have mercy and save us.'' This reference does not treat Eastern
prayer-ropes further; they belong to their own traditions and are not
Rosaries.}

\subsection{The psalter substitution and the lay psalters}

The decisive Western practice is narrower and better documented: the
substitution of counted vocal prayers for the psalter owed to a dead
brother. The obligation itself is visible by 800, when the compact
between St.~Gall and Reichenau bound every priest to a Mass and fifty
psalms for each deceased brother; an Anglo-Saxon charter prescribes that
each monk sing ``two fifties'' (\term{twa fiftig}) for certain
benefactors. As the illiterate lay brothers---the \term{conversi}---became
a distinct body from the choir monks, a simple equivalent had to be
found for men who could not read the psalms at all. The ancient customs
of Cluny, collected by Udalric in 1096, state the substitution in one
clause: whoever is a priest is to sing a Mass, and whoever is not a
priest is to say ``fifty psalms \term{or} as many times the Lord's
Prayer.'' The rule of the Knights Templar, from about 1128, binds the
knights who cannot attend choir to fifty-seven Our Fathers, and to a
hundred a day for a week at the death of a brother.\endnote{Thurston,
``The Rosary'': the St.~Gall--Reichenau compact of 800
(\work{Monumenta Germaniae Historica: Confraternitates}, ed.\ Piper,
140); Kemble, \work{Codex Diplomaticus} I, 290 (\term{twa fiftig}); the
Cluny customary of Udalric, 1096 (``quicunque sacerdos est cantet
missam pro eo, et qui non est sacerdos quinquaginta psalmos aut toties
orationem dominicam,'' PL CXLIX, 776); the Templar rule. This is the
practice the Catechism compresses into one sentence at CCC 2678, on
which see section~\ref{sec:magisterium}.}

Two consequences follow, and both are structural for the Rosary. The
first is the number: the substitute inherits the psalter's hundred and
fifty, and with it the psalter's habit of division into fifties. The
second is the name of the instrument. Strings of pebbles, berries, or
bone discs are attested from the eleventh century---the Countess Godiva
of Coventry (c.~1075) left to a monastery's statue of Our Lady ``the
circlet of precious stones which she had threaded on a cord in order
that by fingering them one after another she might count her prayers
exactly''---and throughout the Middle Ages, across the whole of Europe,
such strings were called \term{paternosters}. Their makers formed
recognized craft guilds: \'Etienne Boileau's \work{Livre des m\'etiers}
records four guilds of \term{paten\^otriers} in Paris in 1268, and
London's Paternoster Row still remembers where the English trade
gathered. The inference Thomas Esser drew in 1897 is difficult to
resist: an appliance persistently called a paternoster was designed to
count Our Fathers, and Godiva's circlet cannot have been meant for Hail
Marys, because in 1075 the Hail Mary was not yet a devotional
formula.\endnote{Thurston, ``The Rosary'': Godiva's bequest (William of
Malmesbury, \work{Gesta Pontificum}, Rolls Series 311); the beads found
in the tomb of St.~Rosalia (1160); the \term{fila de paternoster} and
\term{numeralia de paternoster} of the documents; Boileau's
\work{Livre des m\'etiers}, 1268; Esser's argument of 1897. That the
counting string preceded the Marian psalter is the load-bearing point;
the etymological detail supports it.}

\subsection{The repeated Ave and ``Our Lady's Psalter''}

The Hail Mary entered general devotional use only around the middle of
the twelfth century, and there is ``little or no trace'' of it as an
accepted formula before about 1050; it seems to have grown out of
versicles and responsories of the Little Office of the Blessed Virgin,
two Anglo-Saxon manuscripts of about 1030 already showing its clauses
scattered through the \term{cursus}. Because the words are a greeting
rather than a petition, those who first used them treated them as a
salutation and accompanied each with a bodily act of homage---which is
why repetition, and counted repetition, attached itself to the Ave so
naturally. St.~Aybert in the twelfth century said a hundred and fifty
Hail Marys a day, a hundred with genuflections and fifty with
prostrations; the biographer of St.~Albert (d.~1140) records a hundred
genuflections and fifty prostrations daily, quoting the whole formula as
then said, which ended at ``the fruit of thy womb''; the
\work{Ancren Riwle}, a text older than 1200, directs English anchoresses
to say fifty Aves divided into sets of ten, with prostrations. The
custom of fifty or a hundred and fifty Aves is therefore established in
England and on the Continent before St.~Dominic's birth, and its measure
is openly borrowed: the exercise was called ``Our Lady's Psalter'' for
centuries.\endnote{Thurston, ``Hail Mary,'' \work{The Catholic
Encyclopedia}, vol.~7 (1910),
\url{https://www.newadvent.org/cathen/07110b.htm} (little or no trace
before c.~1050; the Little Office versicles; the British Museum
Anglo-Saxon manuscripts; St.~Aybert; St.~Louis's fifty evening Aves;
the \work{Ancren Riwle} genuflections); Thurston, ``The Rosary''
(St.~Albert with the formula quoted in full; the Eulalia legend, in
which a client accustomed to a hundred and fifty Aves is told to say
fifty and say them more slowly; the \work{Ancren Riwle} sets of ten).
The name \term{rosarium}---a garland of roses---reached the devotion
through the thirteenth-century legend of the monk whose Aves were taken
from his lips as rosebuds and woven into Our Lady's crown; ``Our Lady's
Psalter,'' \term{corona}, \term{chaplet}, and Chaucer's ``pair of
beads'' (bead $=$ prayer) are the older names.}

\subsection{The meditations: the Carthusian contribution}

What makes the Rosary the Rosary is neither the beads nor the number but
the meditation on set mysteries during the repetition---and that is
documented only some two centuries after Dominic's death. Esser's
research, published in \work{Der Katholik} in 1897, established that
attaching meditations to the Aves was the contribution of the Carthusian
Dominic of Prussia (1382--1461).

Dominic was born in Poland, studied at Cracow, drifted, and was
converted at twenty-five through Adolf of Essen, prior of the
charterhouse of St.~Alban near Trier; he entered the Carthusians there
in 1409, served as master of novices at Mainz and as vicar of St.~Alban,
and died in the house he had entered. To the hundred and fifty Aves that
then formed the ``Psalter of Mary'' he had the thought of adding
meditations on the life of Christ and of his Mother. Since the Ave in
his day ended at \term{fructus ventris tui, Jesus}, he simply continued
the sentence, joining to the Holy Name a clause (\term{clausula})
recalling one event of the economy---``\ldots Jesus, \term{quem Angelo
nuntiante de Sancto Spiritu concepisti}''; ``\ldots \term{quo concepto,
in montana ad Elizabeth ivisti}''---and so through a hundred and fifty
clauses. Dominic and Adolf worked to spread the form in the Carthusian
Order and among the laity, and some authors accordingly hold that this
``Psalter'' was one of the original forms from which the present Rosary
developed.\endnote{A.~Mougel, ``Dominic of Prussia,'' \work{The
Catholic Encyclopedia}, vol.~5 (1909),
\url{https://www.newadvent.org/cathen/05112b.htm}, with the biography,
the sample clausulae quoted above, and the caution that ``what has been
improperly called `the Carthusian Rosary' is ascribed to him'';
Thurston, ``The Rosary,'' for Esser's studies in \work{Der Katholik}
(October, November, December 1897). Neither article claims that Dominic
of Prussia invented the counted Aves; the clausulae are a technique of
meditation grafted onto an existing psalter.}

Two things follow. First, the clause technique is not an antiquarian
curiosity: it survives in the regional custom---praised by Paul VI and
by John Paul II---of naming the mystery after the Holy Name in each Hail
Mary, ``done precisely in order to help contemplation and to make the
mind and the voice act in unison'' (MC 46). Second, the fifteen
mysteries as we know them are not yet present. At the close of the
fifteenth century ``the utmost possible variety of methods of
meditating prevailed,'' and the fifteen were ``not uniformly adhered to
even by the Dominicans themselves.''\endnote{MC 46; RVM 33. The
late-fifteenth-century variety: Thurston, ``The Rosary,'' citing
Schmitz, \work{Rosenkranzgebet}, p.~74, and Esser in \work{Der
Katholik} for 1904--6. Neither monograph was examined here.}

\subsection{Alan de la Roche and the confraternities}

The devotion's modern organization---and the Dominic attribution
itself---stem from one Breton Dominican. Alan de la Roche (Alanus de
Rupe) was born about 1428 and died at Zwolle in Holland on 8 September
1475. He entered the order young, studied at Saint-Jacques in Paris, and
from 1459 taught almost without interruption at Paris, Lille, Douai,
Ghent, and Rostock, where he was made Master of Sacred Theology in 1473.
His vision of the restoration of the devotion is assigned to 1460, and
from about that year he preached the re-establishment of ``Our Lady's
Psalter'' with great success throughout northern France, Flanders, and
the Netherlands.\endnote{J.~McNicholas, ``Alanus de Rupe,'' \work{The
Catholic Encyclopedia}, vol.~1 (1907),
\url{https://www.newadvent.org/cathen/01246a.htm}, for the dates,
birthplace (on the testimony of his disciple Cornelius Sneek), teaching
posts, the 1473 doctorate, and the assignment of the vision to 1460.}

The institutional consequence outlasted the man. The Rosary
confraternities organized by Alan and his colleagues at Douai, at
Cologne, and elsewhere ``had great vogue, and led to the printing of
many books, all more or less impregnated with the ideas of Alan.''
Indulgences were granted for the work; and---this is the mechanism that
matters---the documents conceding those indulgences accepted and
repeated the historical data supplied, as was then usual, by the
promoters of the confraternities themselves. A confraternity movement
thus wrote its own foundation narrative into papal instruments, which
later instruments then cited.\endnote{Thurston, ``The Rosary''
(confraternities at Douai and Cologne; the printing; the indulgence
documents repeating promoters' data ``as was natural in that uncritical
age''). The confraternities' own indulgence apparatus is superseded
discipline and is not restated in section~\ref{sec:practice}.}

Alan published nothing in his lifetime. Immediately after his death his
province was ordered to collect his writings; the editions that resulted
``have occasioned much controversy among scholars,'' and his relations
of the visions and sermons of St.~Dominic ``are not to be regarded as
historical.'' The older Dominican judgment, in Qu\'etif and Echard, is
harsher and is not softened here: Alan was earnest and devout, but ``full
of delusions,'' basing his revelations ``on the imaginary testimony of
writers that never existed.''\endnote{McNicholas, ``Alanus de Rupe''
(posthumous collection and editions; relations not historical);
Thurston, ``The Rosary,'' citing Qu\'etif and Echard, \work{Scriptores
Ordinis Praedicatorum} I, 849. The judgment is reported as the
assessment of these authorities; this study did not examine Alan's
printed corpus.}

\subsection{The Hail Mary completed}

The petition ``Holy Mary, Mother of God\ldots'' grew in the same period,
and by an unofficial route. In St.~Louis's time the Ave still ended with
Elizabeth's words; the Holy Name was added later, its introduction
commonly ascribed to Urban IV (1261) on evidence Thurston judges not
clear enough to assert. Because Reformation-era polemic pressed the
charge that a greeting containing no petition is not a prayer, and
because the difficulty had long been felt, those who said their Aves
privately through the fourteenth and fifteenth centuries added closing
clauses of their own---and the surviving verse paraphrases in Italian,
Spanish, German, and Proven\c cal show a general drift toward an appeal
for sinners and for the hour of death. The Bridgettine \work{Myroure of
our Ldy} records both the freedom and its limit: such additions ``may be
saide when folke saye their Aves of theyr own devocyon. But in the
servyce of the chyrche, I trowe it to be moste sewer and moste medeful
to obey the comon use of saying, as the chyrche hath set, without all
such addicions.'' A third part in nearly the modern form appears in a
French \work{Calendar of Shepherds} of 1493 and in Pynson's English
version soon after (``Holy Mary moder of God praye for us synners.
Amen.''); a form identical with the present one but for a single word
(\term{nostrae}) stands printed at the head of a work of Savonarola's
issued in 1495. For liturgical purposes the Ave ended at ``Jesus, Amen''
until the Roman Breviary of 1568 gave the complete prayer official
recognition---and the Catechism of the Council of Trent had already, in
1566, taught that the petition was framed by the Church
herself.\endnote{Thurston, ``Hail Mary,'' for the whole paragraph:
Urban IV and the Holy Name; the Reformers' objection (Latimer,
\work{Works} II, 229--230); the fourteenth- and fifteenth-century
private additions and their vernacular paraphrases; the \work{Myroure
of our Ldy} (Bridgettine nuns of Syon); the Camaldolese and Mercedarian
breviaries c.~1514; the 1493 \work{Calendar of Shepherds} and Pynson's
translation; the 1495 Savonarola printing (British Museum copy) omitting
\term{nostrae}; the 1568 Roman Breviary; the Trent Catechism's
statement that ``most rightly has the Holy Church of God added to this
thanksgiving, petition also and the invocation of the most holy Mother
of God.'' The Trent Catechism is quoted here at second hand from
Thurston's article, not from its own edition.}

\subsection{Pius V and the stabilized form}

St.~Pius V---a Dominican---gave the inherited form its lasting
stability. His bull \work{Consueverunt Romani Pontifices} (17 September
1569) describes the Rosary or ``Psalter of the Blessed Virgin Mary''
exactly: the Blessed Virgin ``venerated by the angelic greeting repeated
one hundred and fifty times, that is, according to the number of the
Davidic Psalter, and by the Lord's Prayer with each decade,'' with
``certain meditations showing forth the entire life of Our Lord Jesus
Christ'' interposed; the bull's operative business is the confirmation
of indults and indulgences for those who pray it. It does not enumerate
fifteen fixed mysteries; Paul VI's measured summary is that Pius V
``explained and in a certain sense established the traditional form of
the Rosary.''\endnote{\work{Consueverunt Romani Pontifices}, English
presentation at \url{https://www.papalencyclicals.net/pius05/p5consue.htm}
(a partial translation whose concluding grant of indults and
indulgences is summarized editorially; translator unnamed); MC 42. The
bull itself repeats the received Dominic attribution (``as is piously
believed''). Lepanto, the feast of 1571--1716, and the calendar are
treated with the current liturgical discipline in
section~\ref{sec:practice}.}

The stabilization was of a form already there. Nothing in the bull
invents the number, the beads, the clauses, or the cycles; it names,
describes, and privileges what confraternities had been praying for a
century---which is exactly what ``in a certain sense established''
concedes and denies at once.

\subsection{The attribution to St.~Dominic: how it rose}

The received account is old, papally repeated, and liturgically
enshrined. The pre-1969 Breviary lesson for the feast defines the
devotion (``a certain form of prayer wherein we say fifteen decades or
tens of Hail Marys with an Our Father between each ten, while at each of
these fifteen decades we recall successively in pious meditation one of
the mysteries of our Redemption'') and then narrates that when the
Albigensian heresy was devastating the country of Toulouse, St.~Dominic
besought Our Lady's help and was instructed by her, ``so tradition
asserts,'' to preach the Rosary; from that time the devotion was ``most
wonderfully published abroad and developed'' (\term{promulgari
augerique coepit}) by St.~Dominic, ``whom different Supreme Pontiffs
have in various past ages of their apostolic letters declared to be the
institutor and author of the same devotion.'' Pius V wrote that Dominic
devised and propagated the method ``as is piously believed''; Leo XIII's
encyclicals assume the institution as historically established; and the
Baltimore Catechism taught American schoolchildren that ``St.~Dominic
taught the use of the Rosary in its present form.''\endnote{Thurston,
``The Rosary,'' for the Breviary lesson, quoted there in full and
reproduced here from that article rather than from a Breviary edition;
\work{Consueverunt} as above; \work{Supremi apostolatus officio} 3
(Dominic ``the first to institute'' the devotion) and 8; \work{Baltimore
Catechism No.~3} (1921), Q.~1081, pp.~240--241 (reception witness). That
many popes have so spoken, Thurston writes, ``is undoubtedly true.''}

The tradition's traceable source is a single preaching campaign. The
Bollandists, trying to trace the story to its origin, found that ``all
the clues converged upon one point, the preaching of the Dominican Alan
de Rupe about the years 1470--75''; he it was ``who first suggested the
idea that the devotion of `Our Lady's Psalter' \ldots\ was instituted or
revived by St.~Dominic.'' The earliest papal bull to carry the story,
Leo X's \work{Pastoris aeterni} of 1520, still hedges it---\term{prout
in historiis legitur}, ``as is read in the histories''---and, in
Thurston's dry summary, ``many of the later popes were less
guarded.''\endnote{Thurston, ``The Rosary'' (Bollandist convergence on
Alan; Leo X's \work{Pastoris aeterni}, 1520). \work{Pastoris aeterni}
was not examined in its own edition here.}

\subsection{The critique, and what it does not establish}

Three findings carry the historical case.

\textbf{Silence where evidence should be.} Of the eight or nine early
Lives of St.~Dominic, not one alludes to the Rosary. The witnesses in his
canonization process are equally reticent; the Balme--Lelaidier
\work{Cartulaire de St.~Dominique} ignores the question; the early
provincial constitutions contain no reference; and among the thousands
of Dominican manuscripts written between 1220 and 1450---sermons,
chronicles, treatises, saints' lives---``no single verifiable passage
has yet been produced which speaks of the Rosary as instituted by
St.~Dominic.'' The charters and deeds of the men's and women's convents
are silent; so are the paintings and sculptures of those two and a half
centuries, the frescoes of Fra Angelico, and the saint's own tomb at
Bologna.

\textbf{The surrendered exhibits.} Each principal proof once offered has
been given up, and mostly by Catholic scholars: the memoirs of
``Luminosi de Aposa,'' who professed to have heard Dominic preach at
Bologna, are a proven forgery; the Muret fresco no longer exists and its
rosary is thought to have been painted in later; the contemporary verses
about a crown of roses rest on a lost manuscript whose only reader read
\term{Dominus}, not \term{Dominicus}; the will of Anthony Sers, leaving a
bequest to a Rosary confraternity at Palencia in 1221, is admitted by
Dominican authorities to be a forgery; and the supposed reference in
Thomas \`a Kempis is a blunder. By Thurston's day the standard German
Catholic reference works had ceased to defend the personal connection at
all.

\textbf{The wrong kind of variety.} A devotion revealed complete and
guarded from the beginning by a great order should not display the
``multitude of conflicting legends'' about its own origin, or the wide
diversity in its manner of recitation, that in fact prevailed down to the
end of the fifteenth century.\endnote{Thurston, ``The Rosary,'' for all
three findings, with the sources he names: Balme and Lelaidier,
\work{Cartulaire de St.~Dominique}; Jean Guiraud's edition of the
\work{Cartulaire de La Prouille} I, cccxxviii; \work{The Month},
February 1901, pp.~179 and 187, and \work{The Irish Rosary}, January
1901, p.~92, for the surrendered exhibits; the \work{Kirchliches
Handlexikon} and Herder's \work{Konversationslexikon} for the change of
tone. None of these was independently examined here; \work{The Month}
articles in particular remain an outstanding item in the terminal
appendix.}

The boundary, then, is exact. The St.~Dominic attribution is received
devotional and papal tradition, and the repetition of that tradition by
popes is itself a fact of reception history---but it is not documentary
evidence that the complete modern form arrived in one revelation.
Conversely, the critique establishes no negative about Dominic's
sanctity, his Marian devotion, or his order's later service to the
Rosary: the same article that dismantles the legend states that the
immense diffusion of the Rosary and its confraternities in modern times,
``and the vast influence it has exercised for good,'' are ``mainly due
to the labours and the prayers of the sons of St.~Dominic.'' The Church
has never defined the Rosary's origin, and the modern papal documents
deliberately state the history in developmental terms---Paul VI himself
noting that the historians' work is ``aimed not at defining in a sort of
archaeological fashion the primitive form of the Rosary but at uncovering
the original inspiration and driving force behind it and its essential
structure.''\endnote{Thurston, ``The Rosary'' (Dominican diffusion; the
order's interest ``only awakened in the last years of the fifteenth
century''); MC 43; RVM 1. Thurston's own answer to the standing
objection of ``vain repetitions'' belongs with the practice section: the
objection ``is felt by none but those who have failed to realize how
entirely the spirit of the exercise lies in the meditation upon the
fundamental mysteries of our faith.''}
